% === Threats to Validity ===
\section{Threats to Validity}
\label{sec:threats}

\noindent\textbf{External validity.} The central limitation is the split between
realism and instrumentation: our telemetry-rich datasets are synthetic, while our
real dataset is text only. As a result, the retrieval and triage claims are
demonstrated on real, engineer-written incidents, but the system's
telemetry-diagnosis behavior---the on-demand evidence gathering that depends on
injected fault structure---is exercised only on synthetic data, and we do not
claim it generalizes to production telemetry; closing that gap is our main future
direction. The real incidents themselves come from open-source Apache projects,
whose public bug tickets may differ from incidents in a proprietary production
setting, and the two synthetic applications, though architecturally distinct, are
demonstration systems rather than a representative sample of all deployments.

\noindent\textbf{Construct validity.} For the real dataset there is no
ground-truth ``which past incident is the answer,'' so matches are inferred from
shared project, component, and symptom overlap rather than human-adjudicated. We
mitigate this with an LLM-as-judge check that separates true matches from random
controls on every dataset (\S\ref{sec:results-rq4}), and we report retrieval under
both a coarse and a strict relation; still, this is a proxy, not a multi-annotator
study, and we release a human-annotation kit to enable one. Triage labels are
likewise derived from each ticket's type and resolution, a reasonable but
imperfect stand-in for whether an engineer truly considered an issue worth acting
on.

\noindent\textbf{Internal validity.} All retrievers enforce time-aware
visibility---only incidents resolved before a window are eligible---so the
evaluation cannot leak future information, and every dataset holds out whole
fault families or project domains so that test items never appear in training.
Experiments use a fixed seed and deterministic decoding for the language models.
We report multi-seed variance for the trained retriever on the synthetic
datasets; on the real dataset, whose single training run is far more costly, we
report a single seed and note this rather than implying broader coverage.

\noindent\textbf{Conclusion validity.} Comparing many systems across datasets and
metrics inflates the chance of spurious findings, so every reported difference is
tested with a paired bootstrap and the family of tests is corrected for false
discovery; we report corrected $q$-values and treat a result as significant only
when its confidence interval excludes zero. The smaller synthetic test sets yield
correspondingly wider intervals, which is why several differences there are
reported as not significant. Finally, we report negative and mixed results
explicitly---fusion's lack of benefit on small corpora, the knowledge graph's
marginal contribution, and triage's collapse on real data---rather than
foregrounding only the favorable comparisons.
